% OWNER: framing
% STATUS: review
% SOURCES: notes/SPIS_TRESCI_PROPOZYCJA_ASP.md, notes/HIPOTEZY_REWIZJA.md, GLOSSARY.md
\chapter{Wprowadzenie}
\label{ch01:top}
\ChapterStatus{draft}{framing}

\section{Wnętrze jako doświadczenie i obraz}
\label{ch01:sec:tlo}

Architektura wnętrz pracuje nie tylko planem i materiałem. Pracuje też \emph{obrazem} proponowanego stanu przestrzeni, czyli wizualizacją, która pośredniczy między zamiarem projektanta a wyobrażeniem użytkownika.
W domowej praktyce ta rola rośnie: o nastroju, materiale i układzie pokoju coraz częściej rozstrzyga ekran, na którym oglądamy moodboardy, rendery, a ostatnio także obrazy generatywne, zanim powstanie jakakolwiek dokumentacja wykonawcza. \textcolor{red}{[ŹRÓDŁO DO UZUPEŁNIENIA]}

Problem nie zaczyna się więc od aplikacji. Zaczyna się od pytania dyscypliny: \emph{jak projektować personalizację koncepcji wnętrza}, skoro doświadczenie przestrzeni jest jednocześnie estetyczne, afektywne i użytkowe, a narzędzie wizualne potrafi w kilka sekund wytworzyć wiarygodny obraz, nie ujawniając przesłanek swojego wyboru?
Rozprawa powstaje w Akademii Sztuki w Szczecinie, w dyscyplinie sztuk plastycznych i konserwacji dzieł sztuki. Interesuje mnie w niej artefakt i proces twórczy: dialog estetyczny przełożony na kontrolowaną syntezę obrazu. Katalog funkcji konsumenckiego oprogramowania nie jest jej przedmiotem.

\section{Personalizacja estetyczna a narzędzia generatywne}
\label{ch01:sec:luka}

Narzędzia nazywane potocznie „AI do wnętrz” oferują zwykle szybką stylizację zdjęcia pomieszczenia: użytkownik wybiera gotowy styl, wpisuje krótki prompt i po kilku sekundach dostaje wynik.
\NeedsCite{przegląd / analiza krytyczna konsumenckich narzędzi AI do wizualizacji wnętrz (peer-review lub raport branżowy z~DOI/ISBN)}
Takie rozwiązanie bywa użyteczne jako produkt. Pomija jednak to, czego potrzebują projektant i badacz: rozróżnienia preferencji \emph{jawnych} i \emph{ukrytych}, nastroju regeneracyjnego przestrzeni, funkcji pokoju, osobowości jako hipotezy projektowej, a przede wszystkim możliwości sprawdzenia, który kanał danych naprawdę kieruje wyborem użytkownika.

Luką, którą wypełniam, nie jest zatem brak kolejnego generatora. Brakuje \textbf{zintegrowanego dialogu badawczo-projektowego}: zestawu instrumentów preferencji osadzonego w doświadczeniu użytkownika, przejrzystej (\emph{explainable}) syntezy briefu, czyli wytycznych przekazywanych generatorowi, oraz ślepej macierzy wariantów, która pozwala empirycznie porównać źródła personalizacji.
Psychologię środowiskową (postrzegana regeneracyjność, biofilia) i pomiar preferencji traktuję jako ramę projektowania, nie jako podręcznik psychometrii.

\section{Problem rozbieżności deklaracji i zachowania}
\label{ch01:sec:problem}

Użytkownicy deklarują styl, paletę i to, co lubią, a w zachowaniu (w szybkim geście wyboru, w czasie reakcji, w wyborze spośród podobnych wariantów) często ujawniają inną logikę gustu.
Kahneman rozróżnia myślenie szybkie i wolne, a Greenwald, McGhee i Schwartz proponują pomiar asocjacji ukrytych za pomocą czasów reakcji \parencite{kahneman2011thinking,greenwald1998iat}.
Na tym tle rozbieżność preferencji jawnych i ukrytych przestaje być zwykłym błędem ankiety. Staje się materiałem do poznania własnego gustu.

Dla projektowania wnętrz wspieranego przez AI wynika stąd konsekwencja praktyczna. System, który ufa wyłącznie deklaracji (dyferencjał semantyczny, lista stylów), może wytworzyć obrazy zgodne z formularzem, a rozmijające się z wyborem, jakiego użytkownik dokona w ślepym porównaniu wariantów.
System, który ufa wyłącznie gestowi, może z kolei zlekceważyć funkcję pokoju i świadome wartości.
Potrzebny jest więc artefakt, który nie ukrywa konfliktu źródeł, lecz go pokazuje i mierzy.

\section{Cel, zakres i pytania badawcze}
\label{ch01:sec:cel}

\paragraph{Cel.}
Chcę zaprojektować, zbudować i ocenić eksperymentalną platformę badawczą IDA (\emph{AI Interior Design Dialogue Research Platform}) jako artefakt Research through Design. Platforma zamienia wieloźródłowy profil preferencji w przejrzysty łańcuch personalizacji koncepcji wnętrz i pozwala w terenie badać rozbieżność preferencji oraz skuteczność kanałów personalizacji \parencite{idaPlatform2026}.

\paragraph{Zakres.}
Zajmuję się koncepcyjnymi wizualizacjami wnętrz mieszkalnych (image-to-image na podstawie zdjęcia pokoju). Dokumentacja budowlana i automatyczny kosztorys leżą poza zakresem.
Produkcyjną generację obrazu w IDA wykonuje Gemini Image w Vertex~AI, a dane badawcze zapisuję w infrastrukturze GCP.
Macierz generacji obejmuje \textbf{sześć} źródeł: \texttt{implicit}, \texttt{explicit}, \texttt{personality}, \texttt{mixed}, \texttt{mixed\_functional}, \texttt{inspiration\_reference}.
Moje eksperymenty w ComfyUI (kontrola geometrii, segmentacja i pokrewne techniki) należą do praktyki pracowni. \textbf{Nie} są środowiskiem uruchomieniowym platformy, ale uzasadniają uproszczenia, które wprowadziłem w potoku terenowym.

\paragraph{Pytania badawcze.}
Oś empiryczną tworzy zestaw \RQ{1}--\RQ{6}, który zastąpił wcześniejszą formułę H1--H5 (szczegóły i plan analiz: rozdz.~\ref{ch06:top}):
\begin{enumerate}
  \item[\RQ{1}] Czy zintegrowany flow badawczy (RtD) umożliwia zebranie wielowymiarowego profilu bez katastrofalnego dropoutu?
  \item[\RQ{2}] Czy i w jakim stopniu preferencje ukryte i jawne się rozchodzą, oraz czy ta rozbieżność wiąże się z zachowaniem w macierzy generacji?
  \item[\RQ{3}] Które źródło w ślepej macierzy sześciu wariantów użytkownicy wybierają i czy wybór zależy od profilu? (rdzeń empiryczny)
  \item[\RQ{4}] Czy cechy i facety Big Five wiążą się z parametrami estetycznymi profilu oraz z preferencją źródła \texttt{personality}? (eksploracyjnie)
  \item[\RQ{5}] Czy odległość PRS \texttt{current}$\rightarrow$\texttt{target} wiąże się z preferencją wariantu \texttt{mixed\_functional}?
  \item[\RQ{6}] Czy uczestnicy postrzegają proces jako zrozumiały i dający sprawczość (agency) przy jednoczesnym zbieraniu danych badawczych?
\end{enumerate}

\section{Teza pracy}
\label{ch01:sec:teza}

\begin{quote}
Platforma IDA, jako artefakt Research through Design, pokazuje, że personalizację koncepcji wnętrz można zaprojektować jako \textbf{dialog estetyczny} (preferencje jawne i ukryte, nastrój, funkcja, osobowość) przełożony na \textbf{kontrolowaną syntezę obrazu}, a nie jako czarną skrzynkę generatora.
Praktyka z modelami dyfuzyjnymi (ComfyUI, kontrola geometrii, segmentacja) i produkcyjny potok IDA (synteza promptu $\rightarrow$ image-to-image w Gemini/Vertex) tworzą razem moją metodę łączenia projektowania wnętrz, psychologii środowiskowej i mediów generatywnych.
\end{quote}

Oś wyników stanowią przede wszystkim \RQ{2} i~\RQ{3}. \RQ{1} i~\RQ{6} wspierają ocenę wykonalności artefaktu, a \RQ{4} i~\RQ{5} rozszerzają eksplorację.
Próba pilotażowa jest mała, dlatego nie obiecuję potwierdzonych efektów. Obiecuję spójny artefakt i ramę pomiaru.

\section{Pozycja metodologiczna: Research through Design}
\label{ch01:sec:rtd}

Przyjmuję paradygmat badań przez projektowanie (Research through Design, RtD).
Frayling rozróżnia badania \emph{into}, \emph{through} i \emph{for} art and design. W wariancie \emph{through} artefakt i proces projektowy stają się sposobem wytwarzania wiedzy \parencite{frayling1993research}.
Zimmerman, Forlizzi i Evenson wprowadzają RtD do badań nad interakcją człowieka z komputerem (HCI) jako metodę, w której zaprojektowany obiekt jest argumentem: ucieleśnia hipotezę o preferowanym stanie świata i podlega krytycznej ocenie \parencite{zimmerman2007rtd}.

Dowód w rozprawie ma trzy warstwy (rozwinięcie w rozdz.~\ref{ch06:top}):
\begin{enumerate}
  \item[(A)] \textbf{Praktyka pracowni}: moje eksperymenty z modelami dyfuzyjnymi w ComfyUI (m.in.\ kontrola geometrii, segmentacja). Jest to wiedza praktyczna autora, a nie silnik produkcyjny IDA.
  \item[(B)] \textbf{Artefakt IDA}: działająca platforma dialogowa z przejrzystą syntezą briefu, macierzą sześciu źródeł i ścieżkami Full i Fast \parencite{idaPlatform2026}.
  \item[(C)] \textbf{Ewaluacja terenowa}: udostępnienie artefaktu uczestnikom online i zbieranie danych pod \RQ{1}--\RQ{6}.
\end{enumerate}
Filozofia produktu (Product-First, Research-Embedded, Minimalist Glass Design) podporządkowuje instrumenty badawcze doświadczeniu użytkownika, a nie odwrotnie.

\section{Wkład własny}
\label{ch01:sec:wklad}

Mój wkład ma trzy rejestry, które się przenikają.

\begin{description}[style=nextline]
  \item[Projektowy]
  Wymyśliłem i zrealizowałem IDA jako dialog estetyczny, a nie kalkulator stylu: awatar pośredniczący w rozmowie, rytm ścieżek Full i Fast, język glass. Dokumentuję uzasadnienia decyzji projektowych (\emph{design rationale}) i świadome odrzucenia.

  \item[Medialny]
  Pracowałem z medium dyfuzyjnym w pracowni (ComfyUI, kontrola struktury przestrzeni, segmentacja) i przeniosłem wybrane wnioski do skalowalnego potoku produkcyjnego (Gemini/Vertex). Zapisuję też, czego świadomie \emph{nie} przeniosłem.

  \item[Badawczy]
  Zoperacjonalizowałem wieloźródłowy profil preferencji i ślepą macierz sześciu źródeł jako eksperyment zawarty w artefakcie. Zastąpiłem zarzuconą formułę H1--H5 ramą \RQ{1}--\RQ{6}. Łączę w ten sposób projektowanie wnętrz, psychologię środowiskową i media generatywne w dyscyplinie sztuk.
\end{description}

\section{Struktura rozprawy}
\label{ch01:sec:struktura}

Rozdział~\ref{ch02:top} buduje ramę przestrzeni, regeneracyjności i biofilii jako kategorii projektowych.
Rozdział~\ref{ch03:top} omawia preferencje jawne i ukryte oraz laddering jako tworzywo dialogu.
Rozdział~\ref{ch04:top} dotyczy medium dyfuzyjnego i mojej praktyki w ComfyUI (warstwa~A).
Rozdział~\ref{ch05:top} wprowadza Big Five i IPIP-NEO-120 jako hipotezę projektową, nie wyrok estetyczny.
Rozdział~\ref{ch06:top} opisuje metodologię RtD, trzy warstwy dowodu i plan analiz \RQ{1}--\RQ{6}.
Rozdział~\ref{ch07:top} przedstawia artefakt IDA: koncepcję, budowę i uzasadnienia projektowe.
Rozdziały~\ref{ch08:top}--\ref{ch10:top} zawierają wyniki pilotażu, dyskusję i wnioski.
